\documentclass[10pt,a4paper]{amsart}
\usepackage[margin=19mm]{geometry}
\usepackage{amsmath,amssymb,mathtools}
\usepackage[scr=rsfs,cal=euler]{mathalfa}
\usepackage{xfrac}
\usepackage[unicode,hidelinks,bookmarks=false]{hyperref}
\newcommand{\ie}{i.e.\ }
\newcommand{\cf}{cf.\ }
\newcommand{\resp}{respectively\ }
\newcommand{\Subsection}{\subsection}
\providecommand{\nobreakdash}{\nobreak\hbox{-}}
\setlength{\emergencystretch}{2em}
\multlinegap0pt
\usepackage{amssymb}
\usepackage{enumerate}
\usepackage[all]{xy}
\usepackage{mathtools}
\usepackage{bm}
\usepackage{float}
\newtheorem{prop}{Proposition}
\newtheorem{theorem}{Theorem}
\usepackage{cleveref}
\crefname{prop}{Proposition}{Propositions}
\crefname{theorem}{Theorem}{Theorems}
\newcommand{\EE}{\mathbb{E}}
\title{Monochromatic components with many edges in random graphs}
\author{Hannah Fox, Sammy Luo}
\date{Source: arXiv:2509.01766v2 (CC BY 4.0)}
\begin{document}
\maketitle

\noindent\emph{Source and licence.} Monochromatic components with many edges in random graphs. Version arXiv:2509.01766v2 (CC BY 4.0), distributed by arXiv under the Creative Commons Attribution 4.0 International licence, http://creativecommons.org/licenses/by/4.0/, as recorded in the arXiv licence metadata for this exact version and shown on the abstract page https://arxiv.org/abs/2509.01766v2. Full licence notice: \texttt{licenses/arxiv-2509.01766-CC-BY-4.0.txt}.\par\smallskip

\noindent\textit{Editorial scope.} Counted unit: the article's theorem thm:multipartite (source0.tex lines 376--381). The article's complete original proof of it is delivered with it (source0.tex lines 382--432, one \texttt{proof} environment of 51 lines). No mathematics is written, repaired or re-derived: the statement and the proof are the article's own lines, in the article's own notation, and no retained line is substituted or re-flowed.  Retained uncounted as the source-local context the proof needs: lines 285--291; lines 312--316; lines 335--345.  Nothing was omitted from any retained span, so the package carries no omission record.  No retained line contains a citation, so the wrapper carries no bibliography and no bibliography entry.  No retained line contains a figure environment, an image-inclusion command or drawing code, so no image is delivered and the package declares no asset dependency.  Editorial formatting only, in addition: an AMS \texttt{amsart} preamble that restates the article's own declarations and macros used by the retained lines, namely the environments \texttt{prop}, \texttt{theorem} on separate counters, together with the standard packages those lines need.  The retained environments print in this document as Proposition 1, Theorem 1 in order of appearance, whereas the article prints the same items with the numbering of its own full document.  The complete native source of the article as distributed in the arXiv e-print for this exact version is bundled unchanged as \texttt{sources/184271/source0.tex}.
\begin{prop}
    \label{prop:density-control}
    For any choice of $p=p_n$ such that $p=\omega(\frac{\log n}{n})$ and any $\varepsilon=\varepsilon_1\in (0,1)$, the following holds as $n$ goes to infinity. 
     In $G=G(n, p)$, with high probability, for every graph $H$ on $V(G)$ satisfying $e(H)\ge \varepsilon \binom{n}{2}$ that is either a clique on a subset of vertices or a complete multipartite graph on $V(G)$,
    \[(1-\varepsilon)\mu \le e(H\cap G) \le (1+\varepsilon)\mu,\]
    where $\mu=\EE[e(H\cap G)]=pe(H)$.
\end{prop}

\begin{prop}
    \label{prop:degree-bound}
    For any $c_0\in (0, 1)$, any choice of $p=p_n$ such that $p=\omega(\frac{\log n}n)$, and any $\varepsilon=\varepsilon_2\in (0,1)$, the following holds as $n$ goes to infinity. 
    In $G := G(n, p)$, with probability $1-o(\frac{1}{n})$, every subgraph $H$ with average degree $\bar d(H)\ge c_0 p(n-1)$ has $v(H)> \bar d(H)(1-\varepsilon)\frac{n}{p(n-1)}$.
\end{prop}

\begin{theorem}
	\label{thm:corollary}
    For any choice of $p=p_n$ such that $p=\omega(\frac{\log n}n)$, any constant $c \in (0, 1)$, and any $\varepsilon = \varepsilon_3\in (0,1)$, the following holds as $n$ goes to infinity. In $G:=G(n,p)$, with probability $1-o(\frac1n)$, the following holds for every complete multipartite graph $M$ on the vertices of $G$:
 
    For every pair of subsets $S, T\subseteq V(G)$ with $e_M(S, T) \ge cn^2$, in the graph $G'=G\cap M$, we have 
    \[
        e_{G'}(S,T)\ge (1-\varepsilon)pe_M(S,T),
    \]
    and furthermore,
    \[e_{G'}(S, T)^2\ge 4(1-\varepsilon)e_{G'}(S) e_{G'}(T).\]
\end{theorem}

\begin{theorem}
	\label{thm:multipartite}
    For any choice of $p=p_n$ such that $p=\omega(\frac{\log n}n)$, any constants $c_1,c_2 \in (0, 1)$, and any $\varepsilon\in (0,1)$, the following holds as $n$ goes to infinity. With probability $1-o(\frac1n)$, in $G:=G(n, p)$, the following holds for every complete multipartite graph $M$ on $V(G)$ with minimum degree at least $c_1 n$:
    
    Let $G':=G\cap M$. Then every subgraph $H$ of $G'$ with $\frac{e(H)}{e(G')}\ge c_2$ has a connected component $H'$ with $e(H')\ge \frac{e(H)^2}{e(G')}(1-\varepsilon)$.
\end{theorem}

\begin{proof}
    Condition on the high probability events of \cref{prop:density-control}, \cref{prop:degree-bound}, and \cref{thm:corollary}, with parameters $c_0=\varepsilon(1-\varepsilon)^2\frac{c_1 c_2}{3}$, $c=\varepsilon(1-\varepsilon)^2\frac{c_1^2 c_2}{3}$, $\varepsilon_1=\varepsilon_2=\varepsilon_3=\varepsilon/3$. Let $H_1,\dots, H_k$ be the connected components of $H$. For any set $S\subset V:=V(G)$, let $f(S)=\frac12e_{G'}(S, V)$. In particular, we have $f(H_1)+\cdots+f(H_k)=f(V)=\frac{1}{2}e_{G'}(V,V)=e(G')$. 
    Then, by the generalized mediant inequality, for some $H'=H_\ell$, we have
    \[\frac{e(H')}{f(V(H'))} \ge \frac{e(H_1)+\cdots+e(H_k)}{f(V(H_1))+\cdots+f(V(H_k))} = \frac{e(H)}{e(G')}\ge c_2.\] 
    If $v(H')\ge \frac{c}{c_1}n$, then since $M$ has minimum degree $\delta(M)\ge c_1 n$, we obtain $e_M(V(H'),V)\ge cn^2$. \cref{thm:corollary} then yields
    \[
        e_{G'}(V(H'),V)^2\ge 4(1-\varepsilon_3)e_{G'}(V(H'))e_{G'}(V)=4(1-\varepsilon_3)e(H')e(G'),
    \]
    so that
    \[
        e(H')\ge \frac{e(H')}{f(V(H'))} f(V(H')) \ge \frac{e(H)}{e(G')}\sqrt{e(H')e(G')(1-\varepsilon_3)},
    \]
    which rearranges to give
	\[e(H')\ge \frac{e(H)^2}{e(G')}(1-\varepsilon_3)\ge \frac{e(H)^2}{e(G')}(1-\varepsilon),\]
    in the high probability event we have conditioned on, as desired. So it suffices to reduce to the case where $v(H')\ge \frac{c}{c_1}n$ for some $H'$ such that $\frac{e(H')}{f(V(H'))}$ is large. 
    
    Consider the vertex set
    \[
        T=\bigcup_{i:|V(H_i)|<\frac{c}{c_1}n}V(H_i).
    \]
    First suppose that $|T|\ge \frac{c}{c_1}n$. In this case, similarly to before, we obtain $e_M(T,V)\ge cn^2$, so that $e_{G'}(T,V))\ge (1-\varepsilon_3)pe_M(T,V)\ge (1-\varepsilon_3)p|T|c_1 n$. Then the average degree of the induced subgraph $H[T]$ is
    \[
        \bar d_{H}(T)=\frac{2e_{H}(T)}{|T|}\ge \frac{e_{H}(T)}{\frac{1}{2}e_{G'}(T,V)}(1-\varepsilon_3)pc_1 n.
    \]
    If $\frac{e_H(T)}{f(T)}\ge \frac{e(H)}{e(G')}\ge c_2$, then
    \[
        \bar d_{H}(T)\ge (1-\varepsilon_3)pc_1 c_2 n \ge c_0 pn,
    \]
    so that $H[T]$ contains a connected component $H_i$ with $\bar d(H_i)\ge c_0 pn$. By \cref{prop:degree-bound}, this implies $v(H_i)\ge \bar d(H_i)(1-\varepsilon_3)\frac{n}{p(n-1)}\ge (1-\varepsilon_3)^2 c_1 c_2 n>\frac{c}{c_1}n$,
    contradicting the definition of $T$. So we must have $\frac{e_H(T)}{f(T)}< \frac{e(H)}{e(G')}$. But then there is some connected component $H'$ with $V(H')\subseteq V\setminus T$ such that
    \[
        \frac{e(H')}{f(V(H'))} \ge \frac{e_{H}(V\setminus T)}{f(V\setminus T)} \ge \frac{e(H)}{e(G')},
    \]
    with $v(H')\ge \frac{c}{c_1}n$, so that we can conclude as before that $e(H')\ge \frac{e(H)^2}{e(G')}(1-\varepsilon)$, as desired.

    It remains to consider the case where $|T|<\frac{c}{c_1}n$. Since $\frac{c}{c_1}\ge (1-\varepsilon_2)c_0$, we must have $\bar d_{H}(T)<(1-\varepsilon_2)^{-1} \frac{c}{c_1}p(n-1)$, or else \cref{prop:degree-bound} leads to a contradiction. So $e(T)<\frac{c^2}{2(1-\varepsilon_2)c_1^2}pn(n-1)$. Then
    \begin{align*}
        \frac{e_{H}(V\setminus T)}{f(V\setminus T)}&\ge \frac{e(H)-e_{H}(T)}{f(V)}\ge \frac{e(H)}{e(G')}-\frac{\frac{c^2}{(1-\varepsilon_2)c_1^2}p\binom{n}{2}}{(1-\varepsilon_1)p\binom{n}{2}}\\
        &\ge \left(1-\frac{c^2}{(1-\varepsilon)^2c_1^2c_2}\right)\frac{e(H)}{e(G')} \ge (1-\varepsilon/3)\frac{e(H)}{e(G')}.
    \end{align*}
    We can then argue as before that for some connected component $H'$ of $H$ with $v(H')\ge \frac{c}{c_1}n$, we have $\frac{e(H')}{f(V(H'))}\ge (1-\varepsilon/3)\frac{e(H)}{e(G')}\ge (1-\varepsilon/3)c_2$. obtaining as before that $e_{G'}(V(H'),V)^2\ge 4(1-\varepsilon_3)e(H')e(G')$, we conclude that
    \[
        e(H')=\frac{e(H')}{f(V(H'))} f(V(H')) \ge(1-\varepsilon/3) \frac{e(H)}{e(G')}\sqrt{e(H')e(G')(1-\varepsilon_3)},
    \]
    which rearranges to give
    \[
        e(H')\ge (1-\varepsilon/3)^2(1-\varepsilon_3)\frac{e(H)^2}{e(G')}\ge (1-\varepsilon)\frac{e(H)^2}{e(G')},
    \]
    as desired. 
    Thus in every case, in the high probability event we have conditioned on, we obtain the desired bound.
\end{proof}

\end{document}
